\chapter*{Apéndices}
\addcontentsline{toc}{chapter}{Apéndices}

Los apéndices amplían, con el nivel de detalle que su extensión requiere, material que se sintetiza en el cuerpo del documento. Se ordenan de lo general a lo específico: primero los perfiles de usuario (personas), luego los instrumentos completos de recolección (encuesta, entrevistas), luego el análisis derivado de esos instrumentos (códigos y triangulación) y, por último, el paquete de reproducibilidad de los resultados del MVP y la planilla del modelo económico-financiero.

\begin{itemize}
    \item[\textbf{A}] \nameref{apx:personas} --- perfiles completos de Santiago (Estudiante), Laura (Docente) y Marcela (Coordinadora de educación a distancia).
    \item[\textbf{B}] \nameref{apx:encuesta} --- instrumento aplicado, diccionario de variables y escalas.
    \item[\textbf{C}] \nameref{apx:entrevistas} --- transcripción íntegra de las tres entrevistas realizadas.
    \item[\textbf{D}] \nameref{apx:cualitativo} --- códigos y temas identificados a partir de las transcripciones.
    \item[\textbf{E}] \nameref{apx:triangulacion} --- triangulación entre encuesta y entrevistas.
    \item[\textbf{F}] \nameref{apx:reproducibilidad} --- repositorio, versión evaluada y comandos para reproducir los resultados.
    \item[\textbf{G}] \nameref{apx:modelo_financiero} --- planilla de cálculo y \textit{scripts} que generan las cifras del modelo de negocio.
\end{itemize}
